\begin{theorembox}{Theorem}
\begin{theorem}[Logarithmic form of inverse hyperbolic sine]
\label{thm:inverse-hyperbolic-sine-log-form}
For every $x\in\mathbb{R}$,
\[
  \operatorname{arcsinh} x=\log\!\left(x+\sqrt{x^2+1}\right).
\]
\end{theorem}
\end{theorembox}
\begin{remark*}[Standard quantified statement]
For every admissible input satisfying the hypotheses of the statement, the
displayed definition or assertion holds in the stated domain.
\end{remark*}

\begin{remark*}[Interpretation]
The statement records the named trigonometric object or identity together with
its intended domain of use.
\end{remark*}

\begin{remark*}[Proof]
Let $t=\operatorname{arcsinh}x$. Then $x=\sinh t=(e^t-e^{-t})/2$. With
$z=e^t>0$, this becomes $z^2-2xz-1=0$. The positive solution is
$z=x+\sqrt{x^2+1}$, so $t=\log z$.
\end{remark*}
\begin{dependencies}
\begin{itemize}
  \item \hyperref[def:inverse-hyperbolic-sine]{Inverse hyperbolic sine}
  \item \hyperref[thm:hyperbolic-exponential-form]{Exponential form of hyperbolic sine and cosine}
\end{itemize}
\end{dependencies}

\begin{theorembox}{Theorem}
\begin{theorem}[Logarithmic form of inverse hyperbolic cosine]
\label{thm:inverse-hyperbolic-cosine-log-form}
For every $x\in[1,\infty)$,
\[
  \operatorname{arccosh} x=\log\!\left(x+\sqrt{x^2-1}\right).
\]
\end{theorem}
\end{theorembox}
\begin{remark*}[Standard quantified statement]
For every admissible input satisfying the hypotheses of the statement, the
displayed definition or assertion holds in the stated domain.
\end{remark*}

\begin{remark*}[Interpretation]
The statement records the named trigonometric object or identity together with
its intended domain of use.
\end{remark*}

\begin{remark*}[Proof]
Let $t=\operatorname{arccosh}x$. Then $t\geq 0$ and
$x=\cosh t=(e^t+e^{-t})/2$. With $z=e^t\geq 1$, this becomes
$z^2-2xz+1=0$. The solution compatible with $z\geq 1$ is
$z=x+\sqrt{x^2-1}$, so $t=\log z$.
\end{remark*}
\begin{dependencies}
\begin{itemize}
  \item \hyperref[def:inverse-hyperbolic-cosine]{Inverse hyperbolic cosine}
  \item \hyperref[thm:hyperbolic-exponential-form]{Exponential form of hyperbolic sine and cosine}
\end{itemize}
\end{dependencies}

\begin{theorembox}{Theorem}
\begin{theorem}[Logarithmic form of inverse hyperbolic tangent]
\label{thm:inverse-hyperbolic-tangent-log-form}
For every $x\in(-1,1)$,
\[
  \operatorname{arctanh} x=\frac{1}{2}\log\!\left(\frac{1+x}{1-x}\right).
\]
\end{theorem}
\end{theorembox}
\begin{remark*}[Standard quantified statement]
For every admissible input satisfying the hypotheses of the statement, the
displayed definition or assertion holds in the stated domain.
\end{remark*}

\begin{remark*}[Interpretation]
The statement records the named trigonometric object or identity together with
its intended domain of use.
\end{remark*}

\begin{remark*}[Proof]
Let $t=\operatorname{arctanh}x$. Then
\[
  x=\tanh t=\frac{e^t-e^{-t}}{e^t+e^{-t}}
  =\frac{e^{2t}-1}{e^{2t}+1}.
\]
Solving gives $e^{2t}=(1+x)/(1-x)$, whose right side is positive because
$-1<x<1$. Taking logarithms gives the formula.
\end{remark*}
\begin{dependencies}
\begin{itemize}
  \item \hyperref[def:inverse-hyperbolic-tangent]{Inverse hyperbolic tangent}
  \item \hyperref[thm:hyperbolic-exponential-form]{Exponential form of hyperbolic sine and cosine}
\end{itemize}
\end{dependencies}

\begin{theorembox}{Theorem}
\begin{theorem}[Logarithmic form of inverse hyperbolic cotangent]
\label{thm:inverse-hyperbolic-cotangent-log-form}
For every $x\in(-\infty,-1)\cup(1,\infty)$,
\[
  \operatorname{arccoth} x=\frac{1}{2}\log\!\left(\frac{x+1}{x-1}\right).
\]
\end{theorem}
\end{theorembox}
\begin{remark*}[Standard quantified statement]
For every admissible input satisfying the hypotheses of the statement, the
displayed definition or assertion holds in the stated domain.
\end{remark*}

\begin{remark*}[Interpretation]
The statement records the named trigonometric object or identity together with
its intended domain of use.
\end{remark*}

\begin{remark*}[Proof]
Let $t=\operatorname{arccoth}x$. Then
\[
  x=\operatorname{coth}t=\frac{e^t+e^{-t}}{e^t-e^{-t}}
  =\frac{e^{2t}+1}{e^{2t}-1}.
\]
Solving gives $e^{2t}=(x+1)/(x-1)$. The quotient is positive exactly on the
stated domain, and taking logarithms gives the formula.
\end{remark*}
\begin{dependencies}
\begin{itemize}
  \item \hyperref[def:inverse-hyperbolic-cotangent]{Inverse hyperbolic cotangent}
  \item \hyperref[thm:hyperbolic-exponential-form]{Exponential form of hyperbolic sine and cosine}
\end{itemize}
\end{dependencies}

\begin{theorembox}{Theorem}
\begin{theorem}[Reciprocal form of inverse hyperbolic secant]
\label{thm:inverse-hyperbolic-secant-reciprocal-form}
For every $x\in(0,1]$,
\[
  \operatorname{arcsech} x=\operatorname{arccosh}\!\left(\frac{1}{x}\right).
\]
\end{theorem}
\end{theorembox}
\begin{remark*}[Standard quantified statement]
For every admissible input satisfying the hypotheses of the statement, the
displayed definition or assertion holds in the stated domain.
\end{remark*}

\begin{remark*}[Interpretation]
The statement records the named trigonometric object or identity together with
its intended domain of use.
\end{remark*}

\begin{remark*}[Proof]
Let $t=\operatorname{arcsech}x$. Then $t\geq 0$ and
$x=\operatorname{sech}t=1/\cosh t$, so $\cosh t=1/x$. Since $0<x\leq 1$,
it follows that $1/x\in[1,\infty)$, and the definition of $\operatorname{arccosh}$
gives the displayed identity.
\end{remark*}
\begin{dependencies}
\begin{itemize}
  \item \hyperref[def:inverse-hyperbolic-secant]{Inverse hyperbolic secant}
  \item \hyperref[def:inverse-hyperbolic-cosine]{Inverse hyperbolic cosine}
\end{itemize}
\end{dependencies}

\begin{theorembox}{Theorem}
\begin{theorem}[Reciprocal form of inverse hyperbolic cosecant]
\label{thm:inverse-hyperbolic-cosecant-reciprocal-form}
For every $x\in\mathbb{R}\setminus\{0\}$,
\[
  \operatorname{arccsch} x=\operatorname{arcsinh}\!\left(\frac{1}{x}\right).
\]
\end{theorem}
\end{theorembox}
\begin{remark*}[Standard quantified statement]
For every admissible input satisfying the hypotheses of the statement, the
displayed definition or assertion holds in the stated domain.
\end{remark*}

\begin{remark*}[Interpretation]
The statement records the named trigonometric object or identity together with
its intended domain of use.
\end{remark*}

\begin{remark*}[Proof]
Let $t=\operatorname{arccsch}x$. Then
$x=\operatorname{csch}t=1/\sinh t$, so $\sinh t=1/x$. The definition of
$\operatorname{arcsinh}$ gives the displayed identity.
\end{remark*}
\begin{dependencies}
\begin{itemize}
  \item \hyperref[def:inverse-hyperbolic-cosecant]{Inverse hyperbolic cosecant}
  \item \hyperref[def:inverse-hyperbolic-sine]{Inverse hyperbolic sine}
\end{itemize}
\end{dependencies}
